% Options for packages loaded elsewhere
\PassOptionsToPackage{unicode}{hyperref}
\PassOptionsToPackage{hyphens}{url}
\documentclass[
  ignorenonframetext,
]{beamer}
\newif\ifbibliography
\usepackage{pgfpages}
\setbeamertemplate{caption}[numbered]
\setbeamertemplate{caption label separator}{: }
\setbeamercolor{caption name}{fg=normal text.fg}
\beamertemplatenavigationsymbolsempty
% remove section numbering
\setbeamertemplate{part page}{
  \centering
  \begin{beamercolorbox}[sep=16pt,center]{part title}
    \usebeamerfont{part title}\insertpart\par
  \end{beamercolorbox}
}
\setbeamertemplate{section page}{
  \centering
  \begin{beamercolorbox}[sep=12pt,center]{section title}
    \usebeamerfont{section title}\insertsection\par
  \end{beamercolorbox}
}
\setbeamertemplate{subsection page}{
  \centering
  \begin{beamercolorbox}[sep=8pt,center]{subsection title}
    \usebeamerfont{subsection title}\insertsubsection\par
  \end{beamercolorbox}
}
% Prevent slide breaks in the middle of a paragraph
\widowpenalties 1 10000
\raggedbottom
\AtBeginPart{
  \frame{\partpage}
}
\AtBeginSection{
  \ifbibliography
  \else
    \frame{\sectionpage}
  \fi
}
\AtBeginSubsection{
  \frame{\subsectionpage}
}
\usepackage{iftex}
\ifPDFTeX
  \usepackage[T1]{fontenc}
  \usepackage[utf8]{inputenc}
  \usepackage{textcomp} % provide euro and other symbols
\else % if luatex or xetex
  \usepackage{unicode-math} % this also loads fontspec
  \defaultfontfeatures{Scale=MatchLowercase}
  \defaultfontfeatures[\rmfamily]{Ligatures=TeX,Scale=1}
\fi
\usepackage{lmodern}
\ifPDFTeX\else
  % xetex/luatex font selection
\fi
% Use upquote if available, for straight quotes in verbatim environments
\IfFileExists{upquote.sty}{\usepackage{upquote}}{}
\IfFileExists{microtype.sty}{% use microtype if available
  \usepackage[]{microtype}
  \UseMicrotypeSet[protrusion]{basicmath} % disable protrusion for tt fonts
}{}
\makeatletter
\@ifundefined{KOMAClassName}{% if non-KOMA class
  \IfFileExists{parskip.sty}{%
    \usepackage{parskip}
  }{% else
    \setlength{\parindent}{0pt}
    \setlength{\parskip}{6pt plus 2pt minus 1pt}}
}{% if KOMA class
  \KOMAoptions{parskip=half}}
\makeatother
\setlength{\emergencystretch}{3em} % prevent overfull lines
\providecommand{\tightlist}{%
  \setlength{\itemsep}{0pt}\setlength{\parskip}{0pt}}
% Slide layout defaults for warning-clean scientific decks.
\usepackage{etoolbox}
\IfFileExists{xurl.sty}{\usepackage{xurl}}{}
\IfFileExists{seqsplit.sty}{\usepackage{seqsplit}}{\newcommand{\seqsplit}[1]{#1}}
\protected\def\breakseq#1{\seqsplit{#1}}
\protected\def\breaktt#1{\begingroup\ttfamily\seqsplit{#1}\endgroup}
\setlength{\emergencystretch}{6em}
\tolerance=5000
\hbadness=10000
\hfuzz=1pt
\setlength{\tabcolsep}{2pt}
\AtBeginEnvironment{longtable}{\tiny\renewcommand{\arraystretch}{0.86}\setlength{\tabcolsep}{1pt}}
\AtBeginEnvironment{tabular}{\tiny\renewcommand{\arraystretch}{0.86}\setlength{\tabcolsep}{1pt}}

% Natbib and cross-reference fallbacks — slides don't load natbib
% or cleveref, but combined-PDF manuscript prose may emit these
% commands. The fallback renders citations as a bracketed key list
% and cross-references as detokenized labels so slides don't fail on
% undefined control sequences or raw underscores. \providecommand is
% a no-op if packages load later.
\providecommand{\citep}[1]{[#1]}
\providecommand{\citet}[1]{#1}
\providecommand{\citealp}[1]{#1}
\providecommand{\citeauthor}[1]{#1}
\providecommand{\citeyear}[1]{#1}
\providecommand{\cref}[1]{\texttt{\detokenize{#1}}}
\providecommand{\Cref}[1]{\texttt{\detokenize{#1}}}
\usepackage{bookmark}
\IfFileExists{xurl.sty}{\usepackage{xurl}}{} % add URL line breaks if available
\urlstyle{same}
% fallback for those not using the hyperref driver hyperxmp:
\makeatletter
\@ifundefined{xmpquote}{\newcommand{\xmpquote}[1]{#1}}{}
\makeatother
\hypersetup{
  hidelinks,
  pdfcreator={LaTeX via pandoc}}

\author{\texorpdfstring{}{}}
\date{}

\begin{document}

\section{Optional Knowledge-Graph
Layer}\label{optional-knowledge-graph-layer}

\begin{frame}[fragile,allowframebreaks]{Optional Knowledge-Graph Layer}
An optional, \textbf{LLM-gated} stage lifts the corpus from
bibliometrics to hypothesis-level evidence. For each record, a local
language model (Ollama, default model \texttt{gemma3:4b}) extracts
structured \emph{assertions} --- each encoding a direction (supports /
contradicts / neutral), a confidence score, and a short natural-language
justification --- against the 6 hypotheses declared in configuration.
Assertions are serialized as RDF-compatible nanopublications
\breakseq{[@kuhn2016decentralized]} and scored by a citation-weighted evidence
function.
\end{frame}

\begin{frame}[allowframebreaks]{Assertion Model}
\protect\phantomsection\label{assertion-model}
Each assertion \(a\) encodes:

\begin{itemize}
\tightlist
\item
  \textbf{Direction}: \(\text{supports}\), \(\text{contradicts}\), or
  \(\text{neutral}\) with respect to a hypothesis \(H\)
\item
  \textbf{Confidence}: a score \(c_a \in [0, 1]\) from the LLM
\item
  \textbf{Citation weight}: \(\log(1 + n_{\text{cites}})\), where
  \(n_{\text{cites}}\) is the citation count of the asserting paper
\end{itemize}

The evidence score for hypothesis \(H\) is:

\[
\text{score}(H) = \frac{\sum_{a \in A(H)^+} c_a \cdot \log(1 + n_{\text{cites}}(a)) -
\sum_{a \in A(H)^-} c_a \cdot \log(1 + n_{\text{cites}}(a))}
{\sum_{a \in A(H)} c_a \cdot \log(1 + n_{\text{cites}}(a))}
\]

where \(A(H)^+\) is the set of supporting assertions and \(A(H)^-\) the
contradicting ones. The score ranges from \(-1\) (all evidence
contradicts) to \(+1\) (all evidence supports).
\end{frame}

\begin{frame}[fragile,allowframebreaks]{Incremental Extraction}
\protect\phantomsection\label{incremental-extraction}
Assertion extraction is \textbf{incremental and resumable}: assertions
are appended to \breaktt{nanopublications.jsonl} at configurable
checkpoint intervals (default: 50 papers). On restart, already-processed
papers are skipped automatically, so a long extraction run that is
interrupted can resume without re-processing. The
\breaktt{-\/-clear-assertions} flag discards previous results for a fresh
start.
\end{frame}

\begin{frame}[allowframebreaks]{Gating and Defaults}
\protect\phantomsection\label{gating-and-defaults}
This stage is entirely optional and never runs in the offline default:
with no language model available it is skipped, and the hypothesis
evidence scores read \emph{pending}. The hypotheses themselves --- their
names and scope --- come from configuration and are reported regardless
of whether the scoring stage has run.

The hypotheses explored in this instance are: H1 Wakefulness Efficacy;
H2 Cognitive Enhancement; H3 Low Abuse Liability; H4 Dopaminergic
Mechanism; H5 Off-label Psychiatric Utility; H6 Tolerability.
\end{frame}

\end{document}
